\subsection{Ein unendliches Gegenbeispiel zur Umkehrung}

Ein Isomorphismus zweier Halbgruppen induziert einen Isomorphismus ihrer
Potenzhalbgruppen. Die Umkehrung gilt im Allgemeinen nicht, wie das
Gegenbeispiel von Mogiljanskaja zeigt.

Die Operationen \(\star_{\mathcal P}\) und \(\diamond_{\mathcal P}\)
bezeichnen dabei die aus \(\star\) und \(\diamond\) abgeleiteten
Mengenprodukte gemäß \FormulaRefAuto{PowerSemigroupProductDef}[def].

\Needspace{20\baselineskip}
\FormulaThmDeltaK[Das Mogiljanskaja-Gegenbeispiel]{%
  \begin{aligned}[t]
    &\exists S,T\,\exists\star,\diamond\;\Bigl(\\[-2pt]
    &\quad\SemigroupAlg{S}{\star}\land
      \SemigroupAlg{T}{\diamond}\land{}\\[-2pt]
    &\quad\neg\Finite{S}\land\neg\Finite{T}\land{}\\[-2pt]
    &\quad\neg\exists f\,
      \SemigroupIso{f}{S,\star}{T,\diamond}\land{}\\[-2pt]
    &\quad\exists F\,\SemigroupIso{F}
      {\PowNE{S},\star_{\mathcal P}}
      {\PowNE{T},\diamond_{\mathcal P}}\Bigr)
  \end{aligned}
}{MogiljanskajaPairCounterexample}{%
  \DeltaRow{Mengen}{S\dsep T}
  \DeltaRow{Funktionen}{\star\dsep\diamond\dsep f\dsep F}
}

Es gibt also zwei unendliche, nicht isomorphe Halbgruppen, deren
Potenzhalbgruppen der nichtleeren Teilmengen isomorph sind. Die
Potenzhalbgruppe bestimmt ihre Grundhalbgruppe somit im Allgemeinen
nicht bis auf Isomorphie. Über den entsprechenden endlichen Fall
entscheidet dieses Gegenbeispiel nicht.

\noindent\emph{Beweis:} Die \MogReadingLink{} enthält die Konstruktion
und den vollständigen erklärenden Beweis. Die \MogProofsLink{} führen
die Definitionen, Hilfsaussagen und Ableitungen aus; ihr
\href{\MogEditionDirectory_B28-mog-proofs.pdf\#mog.statement.MogiljanskajaConstructedPairCounterexample}%
  {konkretes Gegenbeispiel}
liefert die Zeugen des Existenzsatzes.
